% Feladatkiiras
% Csak formai igazítás: a szöveg változatlan, a két nyomtatványoldal egy-egy PDF-oldalra van sűrítve.
\begingroup
\small
\setlength{\parskip}{0pt}
\setlength{\parindent}{0pt}
\noindent
\begin{tabular*}{\textwidth}{@{}l@{\extracolsep{\fill}}r@{}}
\textsc{\textbf{Miskolci Egyetem}} & Sorszám: GEMAN/ISXBF3/2026/BSc \\
Gépészmérnöki és Informatikai Kar & \\
Matematikai Intézet & \\
Analízis Tanszék & \\
\end{tabular*}

\vspace{0.25cm}
\begin{center}
\large\textsc{\textbf{Szakdolgozati Feladat}}
\end{center}
\vspace{0.30cm}

\begin{center}
Gál Olivér István (ISXBF3)\\[1mm]
programtervező informatikus BSc jelölt\\
részére
\end{center}

\vspace{0.25cm}
\noindent\textbf{A szakdolgozat tárgyköre:} WhatsApp Business API-ra épülő kommunikációs platform frontend fejlesztése

\vspace{0.22cm}
\noindent\textbf{A szakdolgozat címe:}

\noindent WaccApp -- WhatsApp Business API-ra épülő kommunikációs platform frontendje

\vspace{0.22cm}
\noindent\textbf{A feladat részletezése:}

\begin{itemize}
    \setlength\itemsep{0pt}
    \setlength\parskip{0pt}
    \setlength\parsep{0pt}
    \setlength\topsep{1pt}
    \setlength\partopsep{0pt}
    \setlength\leftmargini{1.2em}
    \footnotesize
    \item Valósítsa meg a WhatsApp Business API-ra épülő kommunikációs rendszer frontendjét reszponzív webes környezetben!
    \item Készítsen kezelőfelületet az egyedi üzenetküldéshez, a sablonos és szöveges kérdőívek szerkesztéséhez és indításához, valamint a fájl- és médiakezeléshez!
    \item Biztosítsa az időzített kampányok beállításának lehetőségét, a szükséges adatok megadását, ellenőrzését és a kapcsolódó státuszok megjelenítését!
    \item Jelenítse meg vizuálisan a beszélgetéseket, a beérkező válaszokat és a korábbi kommunikációs eseményeket!
    \item Tegye lehetővé az üzenetek és kérdőíves válaszok szűrését, valamint alapvető statisztikai feldolgozását grafikus komponensek és adattáblák segítségével!
    \item Integrálja a frontend működését a backend által biztosított REST API-végpontokkal, és fordítson figyelmet az állapotkezelésre, a hibakezelésre, az alapvető biztonsági szempontokra és a felhasználói élményre!
\end{itemize}

\vspace{0.35cm}

\noindent\textbf{Témavezető(k):} Dr. Hornyák Olivér, egyetemi docens

\vspace{1mm}
\noindent\textbf{Konzulens(ek):} nincs

\vspace{1mm}
\noindent\textbf{A feladat kiadásának ideje:} 2026. \makebox[3cm]{\dotfill}

\vfill

\begin{flushright}
\makebox[6cm]{\dotfill}\\[-1mm]
szakfelelős
\end{flushright}
\endgroup

\newpage
\begingroup
\footnotesize
\setlength{\parskip}{0pt}
\setlength{\parindent}{0pt}
\noindent
\begin{minipage}[t]{0.08\textwidth}
\textbf{1.}
\end{minipage}
\begin{minipage}[t]{0.88\textwidth}
\begin{center}
A szakdolgozat feladat módosítása
\end{center}
\vspace{1mm}
\hspace*{2cm} szükséges.\hfill nem szükséges.

\vspace{0.55cm}
\makebox[5cm]{\dotfill}\hfill\makebox[7cm]{\dotfill}\\[-1mm]
\makebox[5cm]{dátum}\hfill\makebox[7cm]{témavezető(k)}
\end{minipage}

\vspace{0.28cm}

\noindent\textbf{2. A feladat kidolgozását ellenőriztem:}

\vspace{1mm}
\noindent
\begin{tabular*}{\textwidth}{@{}p{0.47\textwidth}@{\extracolsep{\fill}}p{0.47\textwidth}@{}}
témavezető (dátum, aláírás): & konzulens (dátum, aláírás):\\[1.5mm]
\dotfill & \dotfill\\[2.2mm]
\dotfill & \dotfill\\[2.2mm]
\dotfill & \dotfill\\
\end{tabular*}

\vspace{0.30cm}

\noindent\textbf{3. A szakdolgozat beadható:}

\vspace{0.24cm}
\noindent\makebox[5cm]{\dotfill}\hfill\makebox[7cm]{\dotfill}\\[-1mm]
\makebox[5cm]{dátum}\hfill\makebox[7cm]{témavezető(k)}

\vspace{0.30cm}

\noindent\textbf{4. A szakdolgozat} \makebox[3cm]{\dotfill} \textbf{szövegoldalt}

\vspace{0.5mm}
\hspace*{1.2cm}\makebox[3cm]{\dotfill} program protokollt (listát, felhasználói leírást)

\vspace{0.5mm}
\hspace*{1.2cm}\makebox[3cm]{\dotfill} elektronikus adathordozót (részletezve)

\vspace{0.5mm}
\hspace*{1.2cm}\makebox[3cm]{\dotfill}

\vspace{0.5mm}
\hspace*{1.2cm}\makebox[3cm]{\dotfill} egyéb mellékletet (részletezve)

\vspace{0.5mm}
\hspace*{1.2cm}\makebox[3cm]{\dotfill}

\vspace{0.5mm}
\noindent tartalmaz.

\vspace{0.24cm}
\noindent\makebox[5cm]{\dotfill}\hfill\makebox[7cm]{\dotfill}\\[-1mm]
\makebox[5cm]{dátum}\hfill\makebox[7cm]{témavezető(k)}

\vspace{0.30cm}

\noindent
\begin{minipage}[t]{0.08\textwidth}
\textbf{5.}
\end{minipage}
\begin{minipage}[t]{0.88\textwidth}
\begin{center}
A szakdolgozat bírálatra
\end{center}
\vspace{1mm}
\hspace*{2cm} bocsátható.\hfill nem bocsátható.

\vspace{1mm}
\noindent A bíráló neve: \makebox[8cm]{\dotfill}

\vspace{0.24cm}
\makebox[5cm]{\dotfill}\hfill\makebox[7cm]{\dotfill}\\[-1mm]
\makebox[5cm]{dátum}\hfill\makebox[7cm]{szakfelelős}
\end{minipage}

\vspace{0.30cm}

\noindent\textbf{6. A szakdolgozat osztályzata}

\vspace{1mm}
\hspace*{2cm} a témavezető javaslata: \makebox[5cm]{\dotfill}

\vspace{0.5mm}
\hspace*{2cm} a bíráló javaslata: \makebox[5cm]{\dotfill}

\vspace{0.5mm}
\hspace*{2cm} a szakdolgozat végleges eredménye: \makebox[5cm]{\dotfill}

\vspace{0.24cm}
\noindent Miskolc, \makebox[5cm]{\dotfill}

\vfill
\begin{flushright}
\makebox[7cm]{\dotfill}\\[-1mm]
a Záróvizsga Bizottság Elnöke
\end{flushright}
\endgroup
